%% ma: L12-B15-0019
%% lop: 12
%% chuong: 5
%% bai: 15
%% dang: 12.15.15
%% loai: TL
%% muc: VD
%% dap_an: "-"
%% nguon: "(NBV)-10 ĐỀ KIỂM TRA HK2-TOÁN 12-DUNG CHUNG 3 SÁCH.pdf (H02, Đề số 2), Phần 4 Câu 2 (tự luận)"
%% kien_thuc: "Bài 15 (phương trình đường thẳng, vectơ cùng phương)"
%% trang_thai: chinh-thuc
%% ghi_chu: "Câu tự luận. Đáp án gốc khớp (sympy: a = 1, b = 1, k = 2). Đã duyệt 19/10/2026."
\begin{ex}
Trong không gian $Oxyz$, cho hai đường thẳng $d_1\colon\dfrac{x-2}{3}=y+1=\dfrac{z+3}2$ và $d_2\colon\dfrac{x-3}1=\dfrac{y-7}{-2}=\dfrac{z-1}{-1}$. Viết phương trình chính tắc của đường thẳng $d$ cắt $d_1$ và $d_2$ đồng thời đi qua điểm $M(3;10;1)$.
\loigiai{Gọi $A=d\cap d_1$, $B=d\cap d_2$. Khi đó $A(2+3a;-1+a;-3+2a)$, $B(3+b;7-2b;1-b)$ với $a,b\in\mathbb R$.

Ba điểm $A$, $B$, $M$ thẳng hàng nên $\overrightarrow{MA}=k\,\overrightarrow{MB}$, với $\overrightarrow{MA}=(3a-1;a-11;2a-4)$, $\overrightarrow{MB}=(b;-2b-3;-b)$:
\[\begin{cases}3a-1=kb\\a-11=k(-2b-3)\\2a-4=-kb.\end{cases}\]
Cộng phương trình thứ nhất và thứ ba: $5a-5=0$, nên $a=1$. Khi đó $kb=2$ và $-10=-2kb-3k=-4-3k$, suy ra $k=2$, $b=1$.

Vậy $A(5;0;-1)$, $B(4;5;0)$, $\overrightarrow{AB}=(-1;5;1)$. Đường thẳng $d$ đi qua $M(3;10;1)$ có vectơ chỉ phương $(-1;5;1)$:
\[d\colon\dfrac{x-3}{-1}=\dfrac{y-10}5=\dfrac{z-1}1.\]}
\end{ex}
